\documentclass{article}
\usepackage{hyperref}
\usepackage{Style}

\nocite{*} % Comentar si quiero citar
%\addbibresource{bibliografia.bib} % Quitar el comentado si quiero usar bibliografia

\begin{document}

\begin{minipage}{2.5cm}
    \includegraphics[width=2cm]{imagen_puc.jpg}
\end{minipage}
\begin{minipage}{12cm}
    {\sc Pontificia Universidad Católica de Chile\\
    Facultad de Matemáticas\\
    Departamento de Matemática\\
    Profesor: Manuel Cabezas -- Ayudante: Benjamín Mateluna}
\end{minipage}
\vspace{2ex}

{\centerline{\bf Introducción a la Geometría - MAT1304}
\centerline{\bf Ayudantía 21}}
\centerline{\bf 04 de junio de 2026}

\vspace{1cm}
\noindent\textbf{Problema 1.} Muestre que la expresión $x^{2}-4x+y^{2}-2y-4$ describe una 
circunferencia, encuentre su radio y centro.

\vspace{5mm}
\noindent\textbf{Problema 2.} Por el punto fijo $A=(-1,0)$ en la circunferencia $x^{2}+y^{2}=1$ se
traza una cuerda $\overline{AB}$ cualquiera. Hallar la ecuación del lugar geométrico del punto 
medio de $\overline{AB}$.

\vspace{5mm}
\noindent\textbf{Problema 3.} Las parábolas descritas por las ecuaciones $x^{2}-x-4=y$ e 
$y^{2}-5y+1=x$ se intersectan en cuatro puntos. Muestre que tales puntos pertenecen a una 
circunferencia.

\vspace{5mm}
\noindent\textbf{Problema 4.} Desmuestre que la pendiente de la recta tangente a la parábola 
$y=x^{2}$ en el punto $(x_{0},y_{0})$ es igual a $2x_{0}$.

\vspace{1cm}
\noindent\textbf{\underline{Ejercicios Propuestos:}}

\vspace{5mm}
\noindent\textbf{Extra 1.} Describa geométricamente el conjunto de puntos tales que 
$x^{2}-y^{2}=0$ ¿Es una parábola, una elipse o la unión de dos rectas?.

\vspace{5mm}
\noindent\textbf{Extra 2.} Hallar la ecuación polinomial de la curva cuyas ecuaciones paramétricas
son
\begin{equation*}
    x=2+3tan(t) \hhtext{e} y=1+4sec(t)
\end{equation*}

%\printbibliography % Quitar el comentado si quiero usar bibliografia

\end{document}